\section{Strided Writes: Parity, Division, and a Case the Simplifier Does Not Close}\label{sec:pr2}

\subsection{The program}
The test \texttt{forall/pr2} fills an array of even size pairwise with values that depend on a bound
\texttt{MINVAL}, which the loop does not change, and asserts that every cell is at least the bound or zero.
\begin{lstlisting}[language=C]
if (CELLCOUNT % 2 != 0) return 1;
for (i = 1; i <= CELLCOUNT / 2; i++) {
  volArray[i*2 - 2] = (3 >= MINVAL) ? 3 : 0;
  volArray[i*2 - 1] = (1 >= MINVAL) ? 1 : 0;
}
for (i = 0; i < CELLCOUNT; i++)
  __VERIFIER_assert(volArray[i] >= MINVAL || volArray[i] == 0);
\end{lstlisting}
The run is \texttt{-unbounded -teq -simplify -split -simplify -include teq:strides -include
benchmarks/monocera/lemmas.dfy}: the module \texttt{strides} holds the lemmas on strided stores, and
\texttt{lemmas.dfy} the one lemma that is not in the database, \lemma{half\_even}.

\subsection{The trace equation system}
The loop runs from $i = 1$, so its trace of $i$ is $\cnt(1, 1, n)$ and the stores are at $2k$ and $2k + 1$
for the iteration $k = i - 1$. The body has four paths, two conditionals in sequence, of which the
simplifier's context removes the impossible ones (\texttt{MINVAL} $\le 1$ implies \texttt{MINVAL} $\le 3$),
and the merged update has three cases on the bound:
\begin{align}
  V &= \mathsf{trace}(v_0, \mathsf{for}\ \bl{\cnt(0, 1, n)}, \bl{\cnt(1, 1, n)}\ \mathsf{do}\ v, k, i \mapsto {} \nonumber\\
  &\qquad \upd{M \le 1}{v[2k := 3][2k + 1 := 1]}\ \ \upd{1 < M \le 3}{v[2k := 3][2k + 1 := 0]} \nonumber\\
  &\qquad \upd{3 < M}{v[2k := 0][2k + 1 := 0]}) \label{eq:pr1}\\
  &\al{\bl{\cnt(0, 1, n)}}{k \mapsto k < C / 2} \qquad 1 + \intf(n) = i' \qquad \neg(i' \le C / 2) \label{eq:pr2}\\
  &\al{\bl{K'}, \bl{I'}}{k, i \mapsto v'[k] \ge M \vee v'[k] = 0} \qquad i'' = \intf(n') < C \label{eq:pr3}\\
  &\Longrightarrow\ v'[i''] \ge M \vee v'[i''] = 0 \nonumber
\end{align}
with $C$ for \texttt{CELLCOUNT}, $M$ for \texttt{MINVAL}, $v' = \cur(V)$, and $C \bmod 2 = 0$, $1 < C$ from
the harness. The constraint \eqref{eq:pr2} is resolved by \lemma{invs\_counter\_upper} to
$n = 0 \vee (0 < C/2 \wedge \intf(n) - 1 < C/2)$, and with the exit $C/2 < 1 + \intf(n)$ this is the fact
$\intf(n) = C / 2$: the number of iterations is the half of the size, as a quotient.

\subsection{The lemmas}
\begin{align*}
  &\lemma{select\_current\_trace\_stride2\_cases}:\quad R \Leftrightarrow (\neg P \wedge \neg Q)\\
  &\quad\Longrightarrow\ \cur(\mathsf{trace}(a_0, \mathsf{for}\ \bl{\cnt(0, 1, n)}, \bl{\cnt(1, 1, n)}\ \mathsf{do}\ a, k, i \mapsto {}\\
  &\qquad\qquad \upd{P}{a[2k := X_1][2k + 1 := Y_1]}\ \ \upd{Q}{a[2k := X_2][2k + 1 := Y_2]}\\
  &\qquad\qquad \upd{R}{a[2k := X_3][2k + 1 := Y_3]}))[j]\\
  &\qquad = \mathsf{if}\ 0 \le j < 2\,\intf(n)\\
  &\qquad\quad \mathsf{then}\ (\mathsf{if}\ j \bmod 2 = 0\ \mathsf{then}\ (\ite{P}{X_1}{\ite{Q}{X_2}{X_3}})\\
  &\qquad\qquad\qquad\quad \mathsf{else}\ (\ite{P}{Y_1}{\ite{Q}{Y_2}{Y_3}}))\\
  &\qquad\quad \mathsf{else}\ a_0[j]\\[4pt]
  &\lemma{div\_nonpos}:\quad 0 < c\ \Longrightarrow\ (x / c \le 0) \Leftrightarrow (x < c)\\[4pt]
  &\lemma{half\_even}:\quad x \bmod 2 = 0\ \Longrightarrow\ (n = x / 2) \Leftrightarrow (x = 2n)
\end{align*}
The stride lemma reads a cell of the final array: the pairs $2k, 2k + 1$ are disjoint and cover
$0, \ldots, 2n - 1$, outside of which the initial array remains, and the value is chosen by the parity of
the cell and by the guards, which do not depend on the iteration. Its side condition requires the guards to
be exhaustive; without it the update would be unspecified where none holds and the equation unsound (the
guards of a trace need not be complete since \texttt{break}, Section~3 of the paper as implemented). The
lemma fixes two stores and three cases on its left-hand side, as the summaries of the algorithms fix the
fragment structure of one program; \texttt{pr3} to \texttt{pr5}, with $2^s$ combinations of $s$ per-cell
guards, have no lemma.

\subsection{The proof}
Of the 81 rule applications over the goals of \texttt{-split}:
\begin{enumerate}[nosep]
\item \textbf{The reads.} \lemma{select\_current\_trace\_stride2\_cases} fires four times, on the read $v'[k]$
  in the constraint \eqref{eq:pr3} of the checking loop and on $v'[i'']$ in the conclusion (each twice, in
  the two halves of the disjunction), with $P := M \le 1$, $Q := 1 < M \le 3$ and the six values; the
  side condition holds since the simplifier produced the three guards as a case split on $M$. Each read
  is now a conditional of constants below $2\,\intf(n)$, every branch of which satisfies
  $x \ge M \vee x = 0$ under its guard, so the assertion holds for $0 \le i'' < 2\,\intf(n)$.
\item \textbf{The harness.} \texttt{-split} separates the case of no iteration, $C / 2 \le 0$, in which the
  array is the initial one; \lemma{div\_nonpos} turns $C / 2 \le 0$ into $C < 2$, which contradicts
  $1 < C$, three times (one per goal that carries the case).
\item \textbf{The cells beyond the stores.} What remains is $2\,\intf(n) \le i'' < C$ with $\intf(n) = C / 2$
  and $C \bmod 2 = 0$: an empty range, which the simplifier does not see, since it inlines solutions of
  arithmetic equations into polynomials only and $C / 2$ is an atom to it, not a polynomial in $C$. Without
  \lemma{half\_even} this goal remains (the run with the module \texttt{strides} alone leaves exactly it,
  with the constraint of the checking loop fully resolved). \lemma{half\_even} rewrites the fact
  $\intf(n) = C / 2$ under $C \bmod 2 = 0$ into $C = 2\,\intf(n)$, a polynomial equation, which the
  simplifier solves and substitutes: $2\,\intf(n) \le i'' < 2\,\intf(n)$ is false.
\end{enumerate}
The constant strides without cases, \lemma{select\_current\_trace\_stride2} and \lemma{stride3}, close two
further tests in the same way and without the harness,\\
\texttt{array\_doub\_access\_init\_const} (stride 2) and \texttt{array\_tripl\_access\_init\_const} (stride 3).\\
In the first, $N = 100000$ is a literal, the
loop \texttt{i <= N} is a counter of $n$ steps with $\intf(n) = 100001$ decided from the constraint and the
exit, and the read $a'[k]$ for $k \le 200000$ is
\[ \ite{0 \le k < 2\,\intf(n)}{(\ite{k \bmod 2 = 0}{0}{0})}{a_0[k]} \]
(44 rule applications).

\subsection{The idea}
Three obstacles of different kinds in one small program. The stride is a shape of the update, answered by a
lemma that fixes the shape; the quotient in the bound is an obstacle of the simplifier's arithmetic,
answered by a lemma that restates the fact in the form the decision procedure handles (the pattern
``arithmetic in the goal's form'' of Section~\ref{sec:principles:matcher}, here the form the \emph{engine}
wants); and the exhaustiveness of the guards is a soundness premise that the lemma must state, since the
engine does not. \lemma{half\_even} stays outside the database as a workaround for a documented limitation
(TODO.md: solutions are inlined into polynomials only), the proper fix being in the simplifier.
